\noindent\unskip\SourceLeaf{18}Thus
\[\begin{aligned}B&\equiv-NN+1\equiv-(b-c)^2+1\\C&\equiv-NN+b-c\equiv-(b-c)^2+b-c\\D&\equiv-NN-(b-c)\equiv-(b-c)^2-(b-c)\end{aligned}\]
\[\begin{array}{rrrr}
b-c\equiv0&B\equiv1&C\equiv0&D\equiv0\\1&0&0&1\\2&0&1&0.
\end{array}\]
\par\noindent\hspace*{\parindent}Thus, if 3 is a residue modulo $n$, then $N$ is divisible by 3.
\Needspace{10\baselineskip}
\begingroup\small
\begin{center}REMARKS ON THE PRECEDING NOTES CONCERNING CUBIC AND BIQUADRATIC RESIDUES.\end{center}
\noindent\hspace*{\parindent}The note [1] was written by \textsc{Gauss} on the endpaper of the binding of his own copy of the Disquisitiones arithmeticae. The statements under [2] are found in a notebook headed »Uraniae sacrum«, and the two general theorems [3] follow two pages later in the same notebook. The derivations under [4] are taken from several loose slips of paper.

\textsc{Gauss} states at the beginning of his Commentatio prima on the theory of biquadratic residues that in 1805 he began to investigate the theory of cubic and biquadratic residues. In the notebook mentioned above, the derivations given here under [2] and [3] immediately follow notes dated 1804. The »Observatio venustissima« [1], which is less far-reaching than the »Criterien« [2] and [3], can hardly postdate them. We are therefore likely dealing here with \textsc{Gauss}'s earliest investigations and hence with the earliest records of the higher reciprocity laws altogether.

The criteria, like the Observatio, are presumably results of induction. Yet it is very remarkable that the factorization of the quadratic form $aa+27bb$ into its two complex factors, and thus of the number $p$ into its two complex prime divisors, already appears here in the criteria. This had already opened the way for \textsc{Gauss} later to put at least the biquadratic reciprocity law into its simplest form. Incidentally, in certain approaches to proving the general cubic reciprocity law (see the notes immediately following), which in any case date from before 1809, \textsc{Gauss} proceeded, at least at the outset, using rational numbers alone.

The statement made at the beginning of [2] under 3) can be deduced as a special case of the first general theorem under [3]. The latter directly yields the reciprocity law in its most general and simplest form if $p=3m+1$ and $q=3n-1$. For the cubic character of the prime divisor $\pi_1=\dfrac{a+3b\sqrt{-3}}{2}$ with respect to $q$ is given by
\[\left[\frac{\pi_1}{q}\right]\equiv(A-B\sqrt{-3})^{q-1}\quad\;(\mathrm{mod.}\ q)\]
